\section{Logical target: what must be internalized}

\subsection{The ordinary presentation of iPRA}

For the purposes of this report, \iPRA\ is first-order intuitionistic arithmetic with one numerical domain, one function symbol for every primitive-recursive function, the defining equations of those symbols, the usual successor axioms, and induction restricted to quantifier-free formulas.  In the repository, primitive-recursive symbols of arity \(k\) are represented by \code{PR k}; terms over \(n\) free variables by \code{PTm n}; and formulas by \code{Fm PTm n}.  The theory \code{PRAAx} contains projection, composition, and primitive-recursion equations, successor nonzeroness and injectivity, and the quantifier-free induction scheme.

A sequent
\[
 \Gamma\vdash_{\PRAAx} A
\]
is the Lean proposition \code{Prf PRAAx Gamma A}.  The proof object is built from intuitionistic natural deduction, equality, theory axioms, and substitution.  The central structural rule for compilation is simultaneous cut:
\[
 \frac{\Gamma\vdash A \qquad \text{for every }B\in\Gamma,\ \Delta\vdash B}
      {\Delta\vdash A}.
\]
This is exposed as \code{Prf.trans}.  It is exactly what a topological proof compiler needs: a node's local hypotheses can be replaced by the already compiled derivations of its dependencies.

\subsection{The native scalar presentation}

The ConciseIPRA development also contains a much smaller native syntax based on scalar iteration and formulas generated by equations and two binders.  The checked theorems \code{native\_iff\_pra} and \code{pra\_iff\_native} show that the native calculus and the ordinary presentation derive corresponding sequents.  This has two consequences for the framework.

First, it is enough to build object proofs in the ordinary presentation, whose connectives and quantifiers are convenient for application statements.  A finished proof can be exported to the native calculus.  Second, the bridge is not a substitute for the proof: \code{pra\_iff\_native} translates an existing derivation; it does not create one from a semantic theorem.

\subsection{Fixed formulas rather than internal syntax interpretation}

The metatheorems below repeatedly use the phrase ``for every fixed formula.''  This qualification is essential.

Let \(A(\vec x)\) be a fixed quantifier-free formula assembled from equations of primitive-recursive terms by propositional connectives.  At the host level, recurse over the finite syntax tree of \(A\).  For each node, construct a concrete primitive-recursive Boolean term and concrete local derivations.  The resulting Booleanizer is a primitive-recursive function because the formula is fixed.

This is not the same as defining one primitive-recursive function
\[
  \operatorname{Eval}(\ulcorner A\urcorner,\vec x)
\]
that evaluates arbitrary coded primitive-recursive terms.  Such a universal total evaluator cannot itself be primitive recursive.  \CIFER\ always performs syntax recursion before entering the object theory.

\begin{remark}
The host compiler may be general software.  Generality at the host level is harmless because each invocation emits a finite, specialized object program and proof.  The object theory sees only that specialization.
\end{remark}

\subsection{The finite evidence discipline}

Every mathematical notion used by a compiled theorem is divided into four layers.

\begin{definition}[Code]
A code is a natural number or a fixed tuple of natural numbers represented by a primitive-recursive pairing scheme.  All operations on codes used by a theorem have concrete \code{PR} symbols.
\end{definition}

\begin{definition}[Verifier]
A verifier is a fixed primitive-recursive Boolean predicate on codes.  Its object formula is usually an equation \(v(\vec x)=1\), hence quantifier-free.
\end{definition}

\begin{definition}[Evidence]
Evidence is a finite code whose verifier succeeds: a trace, factorization, interval, finite list, compatibility table, equality certificate, or local proof certificate.
\end{definition}

\begin{definition}[Producer]
A producer is a primitive-recursive operation that returns evidence from declared inputs and demands.  When no such operation is available, the desired evidence is represented by a boundary formula rather than by an opaque host constant.
\end{definition}

This discipline separates two proof obligations.  \emph{Verification} proves that a supplied code has the advertised mathematical meaning.  \emph{Production} proves that an appropriate code exists and is obtained by a primitive-recursive operation.  Verification is often elementary even when production is deep.

\subsection{Quantifier-free induction as a trace rule}

Suppose \(S(t,\vec x)\) is the coded state after \(t\) iterations of a fixed primitive-recursive transition, and \(I(t,S(t,\vec x),\vec x)\) is quantifier-free.  The ordinary proof pattern is:
\begin{align*}
 &I(0,S(0,\vec x),\vec x),\\
 &I(t,S(t,\vec x),\vec x)\to I(t+1,S(t+1,\vec x),\vec x),\\
 &I(T,S(T,\vec x),\vec x)\to G(\vec x).
\end{align*}
The first two lines are exactly the premises of \code{PRA.qfInd}; universal elimination specializes the induction theorem at the primitive-recursive fuel \(T\); and the last line may conclude an arbitrary formula \(G\).  Thus the restriction to quantifier-free induction does not require the final theorem to be quantifier-free.  It requires the iterative state invariant to be so.

This observation supports a general design rule:

\begin{quote}
Put all unbounded logical structure outside the loop.  Inside the loop, carry finite Boolean facts, equalities, inequalities, bounded counters, and certificate accumulators.
\end{quote}

\subsection{Bounded quantifiers are finite programs}

A strictly bounded existential
\[
  \exists z<b\;R(z,\vec x)
\]
with decidable \(R\) is represented by a finite disjunction or a Boolean fold.  A bounded universal is represented by a finite conjunction.  The fold is primitive recursive in \(b\), and its reflection theorem is proved by quantifier-free induction on the fold length.  Bounded minimization similarly scans \(0,1,\ldots,b-1\), maintaining a flag and the least witness found.

Accordingly, bounded quantifiers may appear in the human specification without increasing the logical strength of the iterative core.  The compiler expands them into fixed primitive-recursive checkers.

\subsection{Coding tuples, lists, and finite maps}

All finite structures can be reduced to one numerical sort.  For a fixed application, choose one pairing function \(\langle a,b\rangle\) with primitive-recursive projections.  A tuple of fixed dimension is nested pairing.  A list is encoded by length plus entries, or by a bounded sequence code.  A finite Boolean table of size \(m\times n\) is encoded by a bit vector or by a sequence of entries.

The framework does not depend on a particular code, but a formal application must provide the following equations and bounds:
\begin{enumerate}
\item projection after construction;
\item update at the selected index and preservation elsewhere;
\item bounded access default behavior;
\item length and range bounds;
\item extensionality under pointwise equality on the declared range.
\end{enumerate}
These are elementary quantifier-free theorems and should live in one reusable arithmetic interface rather than be reproved by every domain packet.

\subsection{Certified setoids for quotient-like objects}

Literal equality of codes is often too strong.  A certified setoid presentation consists of:
\[
 (C,\;\operatorname{Valid},\;\operatorname{EqCert},\;\operatorname{checkEq}),
\]
where \(C\) is the code type, \(\operatorname{Valid}(x)\) is a Boolean validity condition, and \(\operatorname{checkEq}(x,y,e)=1\) verifies a finite certificate that \(x\) and \(y\) represent the same object.  The application supplies primitive-recursive constructors proving reflexivity, symmetry, and transitivity, together with congruence of operations.

Examples include:
\begin{itemize}
\item rationals as signed numerator/positive denominator pairs with cross-multiplication certificates;
\item projective points with a nonzero scaling relation;
\item localization equality witnessed by a killing exponent;
\item ideal membership witnessed by a finite linear combination;
\item homology equality witnessed by a boundary chain;
\item finite atlases whose overlap equality is witnessed chart by chart.
\end{itemize}

This avoids both extremes: one need not quotient the host type and lose explicit witnesses, and one need not mistake representation equality for mathematical equality.

\subsection{Statement classes}

A mature application should identify its theorem class before proving it.

\begin{description}[style=nextline]
\item[Deterministic theorem.] A primitive-recursive producer returns one code and the verifier proves the exact result.
\item[Bounded relational theorem.] A producer returns a nonempty finite set or relation of valid outputs.
\item[Stable-observable theorem.] The internal relation is nonunique, but every valid output gives the same requested finite observation.
\item[Approximate theorem.] All valid outputs lie in a declared diameter at the requested precision.
\item[Gap-conditional exact theorem.] Approximate evidence plus a positive separation certificate determines one exact label.
\item[Residual conditional theorem.] Downstream derivation is complete, but one or more producer formulas remain in the context.
\end{description}

These classes are not ranked by prestige.  The correct class is the strongest one supported by the available data without importing a hidden choice principle.

\subsection{The formal contract of a node}

A proof-graph node \(v\) has:
\begin{enumerate}
\item a formula \(A_v\);
\item an ordered list of predecessor nodes \(u_1,\ldots,u_k\);
\item a kind: input, local, boundary, or final;
\item for a local or final node, a derivation
\[
 A_{u_1},\ldots,A_{u_k},\Gamma\vdash_{\PRAAx} A_v;
\]
\item optional executable metadata: a producer, verifier, demand transformer, and trace schema.
\end{enumerate}
An input formula is already in \(\Gamma\).  A boundary formula has no local derivation.  The graph is topologically ordered at the proof level.  Precision-demand metadata may contain cycles, but those cycles are solved separately before forward proof composition.

This node contract is the smallest unit at which host generation, finite execution, and object proof can be audited independently.
